\section[陳風]{陳風}

\subsection[宛丘]{宛丘}

子之湯兮，宛丘之上兮，洵有情兮，而無望兮。

坎其擊鼓，宛丘之下，無冬無夏，值其鷺羽。

坎其擊缶，宛丘之道，無冬無夏，值其鷺翿。

\subsection[東門之枌]{東門之枌}

東門之枌，宛丘之栩。子仲之子，婆娑其下。

穀旦于差，南方之原。不績其麻，市也婆娑。

穀旦于逝，越以鬷邁。視爾如荍，貽我握椒。

\subsection[衡門]{衡門}

衡門之下，可以棲遲。泌之洋洋，可以療飢。

豈其食魚，必河之魴。豈其取妻，必齊之姜。

豈其食魚，必河之鯉。豈其取妻，必宋之子。

\subsection[東門之池]{東門之池}

東門之池，可以漚麻。彼美淑姬，可與晤歌。

東門之池，可以漚紵。彼美淑姬，可與晤語。

東門之池，可以漚菅。彼美淑姬，可與晤言。

\subsection[東門之楊]{東門之楊}

東門之楊，其葉牂牂。昏以爲期，明星煌煌。

東門之楊，其葉肺肺。昏以爲期，明星晢晢。

\subsection[墓門]{墓門}

墓門有棘，斧以斯之。夫也不良，國人知之。知而不已，誰昔然矣。

墓門有梅，有鴞萃止。夫也不良，歌以誶之。訊予不顧，顛倒思予。

\subsection[防有鵲巢]{防有鵲巢}

防有鵲巢，邛有旨苕。誰侜予美，心焉忉忉。

中唐有甓，邛有旨鷊。誰侜予美，心焉惕惕。

\subsection[月出]{月出}

月出皎兮，佼人僚兮，舒窈糾兮，勞心悄兮。

月出皓兮，佼人懰兮，舒懮受兮，勞心慅兮。

月出照兮，佼人燎兮，舒夭紹兮，勞心懆兮。

\subsection[株林]{株林}

胡爲乎株林？從夏南。匪適株林，從夏南。

駕我乘馬，説于株野。乘我乘駒，朝食于株。

\subsection[澤陂]{澤陂}

彼澤之陂，有蒲與荷。有美一人，傷如之何。寤寐無爲，涕泗滂沱。

彼澤之陂，有蒲與蕑。有美一人，碩大且卷。寤寐無爲，中心悁悁。

彼澤之陂，有蒲菡萏。有美一人，碩大且儼。寤寐無爲，輾轉伏枕。
